% ============================================================
%  Partie 1 — Modèle de base : MCTNet
% ============================================================

\chapter{Partie 1 — Reproduction de MCTNet}
\label{chap:partie1}

\section{Présentation de MCTNet}

MCTNet (\textit{Multi-stage CNN-Transformer Network}) \cite{wang2024} est une
architecture légère combinant CNN et Transformer pour la classification
pixel par pixel de cultures agricoles. Avec moins de 57 000 paramètres, il
atteint des performances comparables aux modèles bien plus lourds.

\section{Architecture}

\subsection{Vue d'ensemble}

MCTNet enchaîne trois blocs CTFusion suivis d'un Global Max Pooling et d'un
classificateur linéaire :

\begin{equation}
  (B,\,10,\,36)
  \xrightarrow{\text{Stage 1}} (B,\,20,\,18)
  \xrightarrow{\text{Stage 2}} (B,\,40,\,9)
  \xrightarrow{\text{Stage 3}} (B,\,80,\,4)
  \xrightarrow{\text{GMP}} (B,\,80)
  \xrightarrow{\text{Linear}} (B,\,N_c)
\end{equation}

Chaque pixel est décrit par 10 bandes spectrales Sentinel-2 sur 36 pas de
temps. $N_c = 5$ pour l'Arkansas, $N_c = 6$ pour la Californie.

\subsection{Bloc CTFusion}

Le bloc CTFusion est l'unité de base répétée trois fois. Il traite la même
entrée en deux branches parallèles puis fusionne leurs sorties :

\begin{itemize}
  \item \textbf{Branche CNN} (ResBlock1D) — capte les motifs locaux de la
        série temporelle (pics phénologiques, transitions).
  \item \textbf{Branche Transformer} (MHSA, 5 têtes) — modélise les
        dépendances globales entre pas de temps.
\end{itemize}

Les sorties sont concaténées sur la dimension canaux, puis réduites
temporellement par un MaxPool1D (stride 2) :

\begin{equation}
  \text{CTFusion}(x) =
  \text{MaxPool1D}\!\Big(\big[\,\text{CNN}(x),\;\text{Transformer}(x)\,\big]\Big)
  \;:\; (B,C,T) \to (B,\,2C,\,T/2)
\end{equation}

\subsection{ALPE — gestion des données manquantes}

Au Stage 1 uniquement, un module ALPE (\textit{Adaptive Learnable Positional
Encoding}) pondère l'encodage positionnel du Transformer en fonction du masque
de disponibilité (\texttt{input2}). Les pas de temps couverts par des nuages
reçoivent un encodage réduit, limitant leur influence sur l'attention.

\subsection{Bilan paramétrique}

\begin{table}[H]
\centering
\caption{Nombre de paramètres de MCTNet}
\label{tab:params_mctnet}
\begin{tabular}{lrr}
\toprule
\textbf{Composant} & \textbf{Arkansas ($N_c=5$)} & \textbf{California ($N_c=6$)} \\
\midrule
CTFusion Stage 1–3   & \multicolumn{2}{c}{56\,393 (partagé)} \\
Classifier           & 405   & 486 \\
\midrule
\textbf{Total}       & \textbf{56\,798} & \textbf{56\,879} \\
\bottomrule
\end{tabular}
\end{table}

\section{Données}

\subsection{Acquisition et format}

Les données Sentinel-2 ont été acquises via Google Earth Engine sur deux
régions agricoles américaines. Chaque pixel est représenté par :
\begin{itemize}
  \item \texttt{input1} : $(N,\,10,\,36)$ — réflectances spectrales normalisées
  \item \texttt{input2} : $(N,\,36)$ — masque de disponibilité (1 = valide)
  \item \texttt{labels} : $(N,)$ — indice de classe
\end{itemize}

Les labels proviennent du Cropland Data Layer (CDL) de l'USDA.

\subsection{Répartition des échantillons}

\begin{table}[H]
\centering
\caption{Répartition train / val / test}
\label{tab:splits}
\begin{tabular}{lrrr}
\toprule
\textbf{Région} & \textbf{Train} & \textbf{Val} & \textbf{Test} \\
\midrule
Arkansas   & 1\,200 & 300 & 8\,523 \\
California & 1\,440 & 360 & 8\,226 \\
\bottomrule
\end{tabular}
\end{table}

Les pixels sont échantillonnés dans des \textbf{zones dispersées} à travers
la région, afin d'éviter l'autocorrélation spatiale entre les splits.

\section{Protocole d'entraînement}

\begin{table}[H]
\centering
\caption{Hyperparamètres (Table 3 de \cite{wang2024})}
\label{tab:hyperparams}
\begin{tabular}{ll}
\toprule
\textbf{Paramètre}   & \textbf{Valeur} \\
\midrule
Optimiseur           & Adam \\
Taux d'apprentissage & $10^{-3}$ \\
Weight decay         & $10^{-4}$ \\
Batch size           & 32 \\
Époques max          & 200 \\
Têtes d'attention    & 5 \\
Kernel Conv1D        & 3 \\
Dropout              & 0.1 \\
Early stopping       & patience = 20 \\
Scheduler            & ReduceLROnPlateau (factor=0.5) \\
\bottomrule
\end{tabular}
\end{table}

Le meilleur modèle est sélectionné selon le \textbf{F1 macro} sur le jeu de
validation. Les métriques finales (OA, Kappa, F1) sont calculées sur le
test set avec le meilleur checkpoint.

\section{Résultats}

\begin{table}[H]
\centering
\caption{MCTNet — métriques sur le test set}
\label{tab:resultats_mctnet}
\begin{tabular}{lccc}
\toprule
\textbf{Région} & \textbf{OA} & \textbf{Kappa} & \textbf{F1 macro} \\
\midrule
Arkansas    & [résultat] & [résultat] & [résultat] \\
California  & [résultat] & [résultat] & [résultat] \\
\midrule
\multicolumn{4}{l}{\textit{Référence \cite{wang2024}}} \\
Arkansas    & 0.9598 & 0.9421 & --- \\
California  & 0.9502 & 0.9313 & --- \\
\bottomrule
\end{tabular}
\end{table}

% \begin{figure}[H]
%   \centering
%   \includegraphics[width=\textwidth]{figures/courbes_mctnet.png}
%   \caption{Courbes d'apprentissage MCTNet — Arkansas et California.}
%   \label{fig:courbes_mctnet}
% \end{figure}

\cleardoublepage
